% !PRINTED 66
\weberpage{66}
%<<CONT
3. \gesperrt{Zwischen höchstens $m+1$ verwandten $T$-Funktionen $m$ter Ordnung $T, T_1, \ldots, T_m$ besteht eine identische Gleichung von der Form}
\[
LT + L_1T_1 + \cdots + L_mT_m = 0.
\]
Denn besteht diese Relation nicht bereits für $L=0$, so kann man von den in $L_1, L_2, \ldots, L_m$ verfügbaren Konstanten zunächst $\lambda_1, \lambda_2 \ldots, \lambda_m, \mu_1, \mu_2 \ldots, \mu_m$ so bestimmen, daß die sämtlichen Produkte $L_1T_1, L_2T_2, \ldots, L_mT_m$ dieselben Periodizitätsfaktoren erhalten, und dann die $C_1, C_2, \ldots, C_m$ so, daß $m-1$ und folglich alle Nullpunkte der Funktion $L_1T_1+L_2T_2+\cdots+L_mT_m$ mit den Nullpunkten der Funktion $T$ zusammenfallen. Daraus aber folgt, daß das Verhältnis von $L_1T_1+L_2T_2+\cdots+L_mT_m$ zu $T$ gleich einer $T^0$-Funktion $-L$ ist, was zu beweisen ist.

4. \gesperrt{Die Quotienten verwandter $T$-Funktionen können durch Hinzufügung eines Faktors $L$ in doppeltperiodische Funktionen verwandelt werden.}

5. \gesperrt{Setzen wir die Existenz einer $T$-Funktion erster Ordnung $t(u)$ voraus, so läßt sich daraus jede $T$-Funktion $m$ter Ordnung ableiten.}

Es sei nämlich $\gamma$ der Charakter von $t(u)$ und $\alpha_1, \alpha_2, \ldots, \alpha_m$ beliebig gegebene Werte. Es verschwindet dann die Funktion
\[
t(u-\alpha_i+\gamma) = t_i(u), \qquad i = 1, 2 \ldots m
\]
in dem Punkte $\alpha_i$ und allen mit $\alpha_i$ kongruenten Punkten, aber in keinem anderen.

Das Produkt
\[
T(u) = t_1(u)t_2(u) \ldots t_m(u)
\]
ist also eine $T$-Funktion $m$ter Ordnung mit den beliebig gegebenen Nullpunkten $\alpha_1, \alpha_2, \ldots, \alpha_m$.

6. \gesperrt{Hieraus läßt sich leicht beweisen, daß jede doppeltperiodische Funktion, die in einem Periodenparallelogramm nur eine endliche Anzahl von Unstetigkeitspunkten hat, als Quotient zweier $T$-Funktionen darstellbar ist.}

Es sei nämlich $\varphi(u)$ eine doppeltperiodische Funktion mit den Perioden $\omega_1, \omega_2$ und den Unstetigkeitspunkten $\alpha_1, \alpha_2, \ldots, \alpha_m$, die auch teilweise in Unstetigkeitspunkte höherer Ordnung zusammenfallen können.
